\cotefile{157-5}
\scanpage{181}
% manuscript position: top
\mspage{97}
% manuscript position: 1–3
on pourra aussi bien traiter le cas général d'une \uline{$\lim$} de type $J$ quelconque ($J$ \unsure{petit|petite}). Pour ceci, soit
% manuscript position: 4–6
\[
\begin{aligned}
 I=I(J)={}&\text{ens des sous-catégories}\\
 &J'\text{ de }J\text{ telles que}\\
 &\operatorname{card}\operatorname{Fl}(J')\leq\pi,
\end{aligned}
\]
% manuscript position: 7–15
ordonnons $I$ par inclusion. Il est évident que $I\neq\emptyset$ (car $\emptyset\in I$), et \unsure{$=$} que \uline{$I$ est grand devant $\pi$}, car si $(J'_\alpha)_{\alpha\in A}$ avec $\operatorname{card}(A)\leq\pi$ est une famille d'\unsure{élmts|élémts} de $I$, alors la \inserted{sous-}catégorie \inserted{de $J$} engendrée par les $J'_\alpha$ est encore dans $I$ (vérification immédiate), et majore les $J'_\alpha$.
% manuscript position: 15–16
\deleted{On a \unsure{donc} pour $\varphi:J\longrightarrow\operatorname{Ind}'_\pi(C)$ donné, \unsure{considérons}, le foncteur}
% manuscript position: 17
\[
\deleted{I\longrightarrow\operatorname{Ind}'_\pi C\overset{\mathrm{déf}}{=}M}
\]

% manuscript position: 18–21
En fait, la considération de $I(J)$ associé à $J$ va nous permettre de prouver ceci, qui généralise 4.4.13 c) \unsure{rentre|contenu|rentrant} dans le b) :
% manuscript position: 22-margin
% slanted
\margin{\unsure{Comparer|Comp.} avec \unsure{énoncé} \unsure{4.10.3|5.10.3} page 270}

% manuscript position: 22–23
\uline{Corollaire 4.4.14} \quad Soit $M$ une catégorie, $\pi$ un cardinal infini tel que a) $M$ est
% manuscript position: foot
% archive sheet number: 180
\scanpage{182}
% manuscript position: top
\mspage{98}
% manuscript position: 1–3
stable par $\underline{\lim}$ de cardinal $\leq\pi$ et b) $M$ est stable par $\varinjlim$ \inserted{sur des ens ord $I$} grands devant $\pi$. Alors $M$ est stable par petites $\varinjlim$.
% manuscript position: left/1–9
% margin note written sideways
\margin{\uline{Cor. 4.4.14.1} $M$ étant comme dans 4.4.14, et $f:M\longrightarrow\unsure{\mathcal E|E}$ un foncteur qui commute aux \inserted{petites} $\varinjlim$ de cardinal $\leq\pi$ et aux $\varinjlim$ \deleted{\unsure{qq}} \unsure{grandes devant} $\pi$, alors $f$ commute aux petites $\varinjlim$.}

% manuscript position: 4–7
\deleted{\uline{NB} On sait que ces hypothèses sont vérifiées pour $\operatorname{Ind}_\pi C$, sous les conditions 4.4.13, grâce à 4.4.13 b) déjà établi et \deleted{\unsure{au}} lemme 4.4.8.1}

% manuscript position: 8–9
En effet, si $\varphi:J\longrightarrow M$ est donné, on en \unsure{conclut|construit}
% manuscript position: 10
\[
\psi:I=I(J)\longrightarrow M
\]
% manuscript position: 11
par
% manuscript position: 12
\[
\psi(J')=\varinjlim_{J'}\varphi\mid J',
\]
% manuscript position: 13–18
qui existe par hypothèse sur $M$. Il est alors immédiat que $\varinjlim_J\varphi$ existe sss $\varinjlim_I\psi$ existe, or cette existence est assurée du fait que $I$ est grand devant $\pi$, et l'hypothèse sur $M$.

% manuscript position: 19–22
Pour achever de prouver le \unsure{th.} 4.4.13, \deleted{\unsure{on}} est donc \unsure{amené|ramené} au cas particulier de limites inductives $\varinjlim_J\varphi(j)$, qd $\varphi:J\deleted{\to}\longrightarrow\operatorname{Ind}'_\pi C$ est donné, avec $J$ \unsure{ens|une}
% manuscript position: foot
% archive sheet number: 181
\scanpage{183}
% manuscript position: top
\mspage{99}
% manuscript position: 1–2
ordonné grand devant $\pi$. Mais ceci va résulter du résultat plus précis :

% manuscript position: 3–7
\uline{Corollaire} 4.4.15. Soit $C$ une petite catégorie. Alors \inserted{la sous-catégorie strictement pleine} $\operatorname{Ind}'_\pi(C)$ \deleted{\unsure{est une sous-}} de $\widehat C$ est stable \inserted{dans $\widehat C$} par $\varinjlim$ grandes devant $\pi$, \deleted{\unsure{en fait}} \deleted{et} \inserted{en d'autres termes} ces $\varinjlim$ \deleted{\unsure{nt}}\unsure{existent|existe} \uline{dans $\operatorname{Ind}'_\pi(C)$}, et le foncteur d'inclusion y \unsure{commute}.
% manuscript position: left/4–7
% slanted
\margin{\unsure{Comparer|Comp.} \unsure{à|avec} 4.7.3.1 \unsure{qui pourrait|qu'il y aurait} \unsure{pu|se} s'expliciter \unsure{ici}.}

% manuscript position: 8–13
\uline{Dém} Soit $I$ ens. ordonné, grand devant $\pi$, et \deleted{$\unsure{(F_i)}$} $i\longmapsto F_i$ un foncteur $I\to\operatorname{Ind}'_\pi(C)$, $F=\varinjlim_I F_i$ dans $\widehat C$, prouvons que $F\in\operatorname{Ob}\operatorname{Ind}'_\pi(C)$, i.e. (4.4.3) que $C_{/F}$ est grd devant $\pi$. J'utilise le critère b) de 4.4.4, il faut vérifier $\mathrm{PF}_{\pi}1$, $\mathrm{PF}_{\pi}2$. Soit
% manuscript position: 14
\[
x_\alpha\xrightarrow{u_\alpha}F=\varinjlim F_i
\]
% manuscript position: 15–23
des objets de $C_{/F}$ indexés par $A$ avec $\operatorname{card} A\leq\pi$. Chaque $u_\alpha$ provient de $x_\alpha\to F_{i_\alpha}$, et en prenant un majorant commun des $i_\alpha$, \inserted{($I$ grd devant $\pi$)} on trouve un $i\in I$ et des $v_\alpha:x_\alpha\to F_i$. Comme $F_i$ est \deleted{\unsure{dans}} $\operatorname{Ind}'_\pi(C)$ i.e. $C_{/F_i}$ grd devant $\pi$, considérant les $v_\alpha$ comme des objets de $C_{/F_i}$, ils sont majorés par un objet $x\longrightarrow F_i$ de $F_i$. On a donc
% manuscript position: 24–25
\diagram{\begin{tikzcd}[ampersand replacement=\&]
x_\alpha \arrow[r] \arrow[d] \& F_{i_\alpha} \arrow[d] \arrow[dr,dashed] \& {}\\
x \arrow[r] \& F_i \arrow[r] \& F\text{,}
\end{tikzcd}}
% manuscript position: 26
qui montre que les $x_\alpha\longrightarrow F$ sont majorés par $x\to F$,
% manuscript position: foot
% archive sheet number: 182
\scanpage{184}
% manuscript position: top
\mspage{100}
% manuscript position: 1–5
ce qui prouve $\mathrm{PF}_\pi1$. Prouvons $\mathrm{PF}_\pi2$, i.e. soient $x\xrightarrow{u}F$, $y\xrightarrow{v}F$ deux objets de $C_{/F}$, et $(\deleted{\unsure{u}}f_\alpha)_{\alpha\in A}$ une famille de cardinal $\leq\pi$ de flèches de $u$ dans $v$.
% manuscript position: left/4–5
\diagram{\begin{tikzcd}[ampersand replacement=\&,column sep=small,row sep=small]
x\arrow[rr,"f_\alpha"]\arrow[dr,"u"'] \& {} \& y\arrow[dl,"v"]\\
{} \& F \& {}
\end{tikzcd}}
% manuscript position: 5–7
On sait que $u,v$ proviennent de $\deleted{\unsure{f}}u_i:x\longrightarrow F_i$, $v_i:y\longrightarrow F_i$ pour $i$ assez grand.
% manuscript position: left/7–9
\diagram{\begin{tikzcd}[ampersand replacement=\&,column sep=small,row sep=small]
x\arrow[rr,"f_\alpha"]\arrow[dr,"u_i"'] \& {} \& y\arrow[dl,"v_i"]\\
{} \& F_i \& {}
\end{tikzcd}}
% manuscript position: 7–16
Pour $\unsure{\forall|\mathrm{tt}}\,\alpha$, on peut choisir $i_\alpha\geq i$ tel que $v_if_\alpha$ et $u_i$ deviennent égaux en composant avec $F_i\longrightarrow F_{i_\alpha}$, et comme $\operatorname{card} A\leq\pi$ et $I$ grd devant $\pi$, il \deleted{\unsure{faut}} \inserted{existe} un majorant $i'$ des $i_\alpha$. Au total, \deleted{\unsure{les}} \unsure{OPS|O.P.S.} (remplaçant $i$ par $i'$) que $v_i f_\alpha=u_i$ pour tout $\alpha$. Alors $u_i,v_i$ \deleted{\unsure{étant}} considérés comme objets de $C_{/F_i}$, les $\unsure{u_\alpha|f_\alpha}$ sont des flèches de $C_{/F_i}$ de $u_i$ dans $v_i$.
% manuscript position: right/16
\margin{$\unsure{=}$}

% manuscript position: 17–19
Comme $C_{/F_i}$ est grd devant $\pi$ (par l'hypothèse $F_i\in\operatorname{Ob}\operatorname{Ind}'_\pi(C)$), on trouve qu'il existe un égalisateur $(y,v_i)\xrightarrow{\unsure{g_i|g}}(z,w_i)$.
% manuscript position: left/20–22
\diagram{\begin{tikzcd}[ampersand replacement=\&,column sep=small,row sep=small]
x\arrow[r,shift left=1,"{(f_\alpha)}"]\arrow[r]\arrow[r,shift right=1]\arrow[dr,"u_i"'] \& y\arrow[r,"g"]\arrow[d,"v_i"] \& z\arrow[dl,"w_i"]\\
{} \& F_i \& {}
\end{tikzcd}}
% manuscript position: 20–27
Alors l'image de $g$ dans $C_{/F}$ (via $F_i\longrightarrow F$) égalise les flèches $f_\alpha$, en tant que flèches de $C_{/F}$, ce qui prouve $\mathrm{PF}_\pi2$ et achève la démonstration de 4.4.15, et par là aussi celle du th. 4.4.13, et du lemme 4.3. Par là est aussi achevée la dém. de la première assertion du lemme
% manuscript position: foot
% archive sheet number: 183
\scanpage{185}
% manuscript position: top
\mspage{101}
% manuscript position: 1–4
4.3.7 (p. 77), ce qui achève la démonstration du th. 4.3,\unsure{5|6} (p. 70-71). Il reste à établir la propriété universelle énoncée au 4.3.7, que je vais rappeler
% manuscript position: left/5–13
% slanted
\margin{Devrait être incorporé en partie d) au th. 4.4.13, \unsure{ceci étant|car c'est} \unsure{un complément|conceptuellement} \unsure{important}, et pas tellement \unsure{qu'évident|évident} finalement.}

% manuscript position: 5–9
\uline{Corollaire} 4.4.16. Soit $C$ une catégorie équivalente \unsure{à une|une} petite catégorie, stable par $\varinjlim$ de cardinal $\leq\pi$, où $\pi$ est un cardinal infini donné, et considérons le foncteur d'inclusion canonique
% manuscript position: 10
\[
 C\xrightarrow{\alpha}\operatorname{Ind}_\pi(C).\tag*{(4.4.16.1)}
\]

% manuscript position: 11–16
Ce foncteur est 2-universel pour les foncteurs $f:C\longrightarrow\unsure{\mathcal E|E}$, où $\mathcal E$ est une \deleted{\unsure{$\mathcal U$-}}\inserted{$\mathcal U$-}catégorie stable par \inserted{petites} $\varinjlim$ \deleted{de cardinal $\leq\pi$}, et $f$ un foncteur commutant aux limites inductives. En d'autres termes, on a

% manuscript position: 17–20
1°) $\alpha$ possède les propriétés précédentes, i.e. $\unsure{\operatorname{Ind}_\pi(C)|\operatorname{Ind}'_\pi(C)}$ est stable par $\varinjlim$, \deleted{de cardinal $\leq\pi$} et le foncteur $\alpha$ commute \unsure{aux} $\varinjlim$ de \deleted{\unsure{de}} cardinal $\leq\pi$ ;

% manuscript position: 21–23
2°) Pour \deleted{\unsure{toute}} \inserted{$\mathcal U$-}catégorie $\mathcal E$ stable par petites limites inductives, le foncteur $g\longmapsto g\circ\alpha$
% manuscript position: 24
\[
\underline{\mathrm{Hom}}_{!}(\operatorname{Ind}_\pi C,\mathcal E)\xrightarrow{\varphi}\underline{\mathrm{Hom}}_{!\pi}(C,\mathcal E)\tag*{(4.4.16.\unsure{1|2})}
\]
% manuscript position: 25–27
est une équivalence de catégories, où l'indice $!$ (resp $!\pi$) indique la sous-catégorie \inserted{str.} pleine de la catégorie $\underline{\mathrm{Hom}}$
% manuscript position: foot
% archive sheet number: 184
\scanpage{186}
% manuscript position: top
\mspage{102}
% manuscript position: 1–3
envisagée, formée des foncteurs qui commutent aux petites $\underline{\lim}$ (resp à celles de cardinal $\leq\pi$).

% manuscript position: 4–16
On peut donc dire aussi que $\operatorname{Ind}_\pi(C)$ joue le rôle d'une ``enveloppe'' stable par \inserted{petites} $\varinjlim$ de la catégorie \deleted{\unsure{$M$}} $C$, quand les $\varinjlim$ \inserted{de cardinal $\leq\pi$, lesquelles} \deleted{qui} existent déjà dans $C$, sont respectées par le plongement de $C$ dans son enveloppe. C'est un résultat \unsure{nullement|réellement} tautologique qu'il suffit pour cela d'ajouter à $C$ des $\varinjlim$ \uline{grandes devant} $\pi$ d'objets de $\unsure{C|c}$, et que \unsure{plus} précisément, la catégorie enveloppe est réalisée par $\operatorname{Ind}_\pi C$ (\unsure{ou}, au choix et moins pléthoriquement, par $\operatorname{Ind}'_\pi C\subset\widehat C$).

% manuscript position: 17–23
\uline{Démonstration} \inserted{Le 1°) est contenu dans 4.4.13 \unsure{a)|b)} et c). \unsure{Prouvons donc|Prouvons le} 2°) \deleted{\unsure{soit|on}}.} Il est évident que $\alpha$\deleted{\unsure{$(C)$}} est un foncteur pleinement fidèle, dont l'image essentielle engendre strictement $M\equiv\operatorname{Ind}_\pi C$, d'où résulte que le foncteur (4.4.16.2) est pl fid, étant induit par le foncteur pleinement fidèle
% manuscript position: 24
\[
\alpha_{\unsure{\mathcal E|E}}^{*}:\underline{\mathrm{Hom}}(M,\mathcal E)\longrightarrow\underline{\mathrm{Hom}}(\deleted{\unsure{M}}C,\mathcal E)
\]
% manuscript position: foot
% archive sheet number: 185
\scanpage{187}
% manuscript position: top
\mspage{103}
% manuscript position: 1–2
sur les sous-catégories pleines envisagées. Il reste donc à montrer que (\unsure{4.4.15.2|4.4.16.2}) est ess. surjectif.

% manuscript position: 3–4
Soit donc $f$\inserted{$:C\longrightarrow\unsure{\mathcal E|E}$} dans $\underline{\mathrm{Hom}}_{!\pi}(C,\mathcal E)$, considérons son extension canonique
% manuscript position: 5–6
\[
\left\{\begin{aligned}
 Lf &: \operatorname{Ind}(C)\longrightarrow\mathcal E\\
 Lf\big((\unsure{X_i|X'})_{i\in I}\big)&=\varinjlim_{i\in I} f(\unsure{X'|X_i})
\end{aligned}\right.\tag*{(\unsure{4.4.13.3|4)4.13.3})}
\]
% manuscript position: 7
(cf SGA 4 I 8.7.3), et soit
% manuscript position: 8
\[
\unsure{\psi|\varphi}(f)=Lf\mid\operatorname{Ind}_\pi(C)\tag*{(4.4.13.4)}
\]

% manuscript position: 9–10
Rappelons que $Lf$ commute aux \inserted{petites} $\varinjlim$ filtrantes, mais comme le foncteur d'inclusion
% manuscript position: 11
\[
\operatorname{Ind}_\pi C\longrightarrow\operatorname{Ind}(C)\tag*{(\unsure{4.4.16.5|4.4.16.3}\text{\deleted{\unsure{3}}})}
\]
% manuscript position: 12–25
ne commute pas à \deleted{\unsure{petites} ces} $\varinjlim$, \deleted{\unsure{nous}} il n'est pas clair a priori que $\psi(f)$ commute aux dites limites, ni a fortiori qu'il commute aux petites $\varinjlim$ sans plus. Mais on a vu que \deleted{$\operatorname{Ind}_\pi C$} \unsure{ce|ces} commute aux $\varinjlim$ \uline{grandes devant $\pi$}, ce qui est contenu dans le lemme 4.4.15 (car que les limites inductives \uline{filtrantes} de $\operatorname{Ind}(C)$ sont les limites \uline{dans $\widehat C$}) : dire que $\operatorname{Ind}'_\pi(C)\mathrel{\unsure{\relbar|\to}}\operatorname{Ind}'_{\unsure{\aleph_0|\pi}}(C)$ commute à un certain type de $\varinjlim$ filtrantes, revient à dire qu'il en est ainsi du composé avec $\operatorname{Ind}'_\pi(C)\mathrel{\lhook\joinrel\relbar}\widehat C$, ce qui est établi par 4.4.15. De ceci résulte que $\psi(f)$ \uline{commute aux} $\varinjlim$ \uline{grdes devant $\pi$}. Pour établir que $g=\psi(f)$ commute
% manuscript position: foot
% archive sheet number: 186
\scanpage{188}
% manuscript position: top
\mspage{104}
% manuscript position: 1–7
aux petites $\varinjlim$, il reste donc à voir qu'il commute aux $\varinjlim$ de cardinal $\leq\pi$ ; en vertu de 4.4.14.1. (marge page 98). Pour l'établir, on utilise le th. 4.4.13. b), en représentant un foncteur $\varphi:J\longrightarrow\operatorname{Ind}_\pi(C)$ \inserted{(où $\operatorname{card}\operatorname{Fl} J\leq\pi$)} sous forme indicielle par
% manuscript position: 8
\[
\Phi:I\times J\longrightarrow C,\tag*{(4.4.16.6)}
\]
% manuscript position: 9–10
$I$ étant un ens. ordonné grand devant $\pi$. On a donc
% manuscript position: 11
\[
\varphi(j)=\varinjlim_{i\in I}{}_{\operatorname{Ind}_\pi C}\Phi(i,j),
\]
% manuscript position: 12–13
et comme $g=\psi(f)$ commute aux $\varinjlim$ de type $I$ par ce qu'on a vu, on aura
% manuscript position: 14
\[
\deleted{g(\varinjlim{}_{\operatorname{Ind}_\pi C}\varphi(j))}
\]
% manuscript position: 15
\[
g(\varphi(j))=\varinjlim_{i\in I}{}_{(\unsure{\mathcal E|E})}g(\Phi(i,j))
\]
% manuscript position: 16
d'où
% manuscript position: 17
\[
\varinjlim_{j\in J}{}_{(\mathcal E)}g(\varphi(j))=\varinjlim_{j\in J}\varinjlim_{i\in I}
\underbrace{g(\Phi(i,j))}_{\substack{\Vert\\ f(\Phi(i,j))}},\tag*{(\(\ast\))}
\]
% manuscript position: right/17
\inserted{(limites \unsure{dans|ds} $\mathcal E$)}

% manuscript position: 18
D'autre part on a
% manuscript position: 19–20
\[
\begin{aligned}
\varinjlim_{j\in J}{}_{\operatorname{Ind}_\pi}\varphi(j)
 &\mathrel{\unsure{=|\simeq}}\varinjlim_{\unsure{J|j}\in J}\varinjlim_{i\in I}\Phi(i,j)\\
 &\simeq\varinjlim_{i\in I}\underbrace{\varinjlim_{j\in J}\Phi(i,j)}_{\substack{\text{\unsure{peut être|peut-être} calculé}\\\text{dans }C,\text{ soit}\\\Psi(i)\in\operatorname{Ob} C}},
\end{aligned}
\]
% manuscript position: 21
donc
% manuscript position: 22
\[
\varinjlim_{j\in J}{}_{\operatorname{Ind}_\pi C}\varphi(j)\mathrel{\unsure{=|\simeq}}\varinjlim_{i\in I}\unsure{\Psi|\Psi'}(i),\tag*{(\(\ast\ast\))}
\]
% manuscript position: foot
% archive sheet number: 187
\scanpage{189}
% manuscript position: top
\mspage{105}
% manuscript position: 1–3
d'où en appliquant $g$, qui commute aux $\varinjlim$ de type $I$, et se rappelant maintenant que $g\mid C\simeq f$. \deleted{et que $f$}
% manuscript position: 4
\[
g\big(\varinjlim_J{}_{\operatorname{Ind}_\pi C}\varphi(j)\big)
\simeq\varinjlim_{i\in I}f(\Psi(i)).\tag*{(\(\ast\ast\ast\))}
\]

% manuscript position: 5
On a d'autre part
% manuscript position: 6
\[
\Psi(i)\overset{\mathrm{déf}}{=}\varinjlim_{j\in J}\Phi(i,j)
\qquad\varinjlim\ \text{dans }C
\]
% manuscript position: 7–8
et comme $f$ commute aux $\varinjlim$ de type $J$ (car $\operatorname{card}\operatorname{Fl} J\leq\pi$), on en conclut
% manuscript position: 9
\[
f(\Psi(i))\simeq\varinjlim_{j\in J}f(\Phi(i,j)),
\]
% manuscript position: 10
donc par $(\ast\ast\ast)$
% manuscript position: 11
\[
g\big(\varinjlim_{j\in J}\varphi(j)\big)\simeq\varinjlim_{i\in I}\varinjlim_{j\in J}f(\Phi(i,j)).\tag*{(\(\ast\ast\ast\ast\))}
\]

% manuscript position: 12–14
En comparant avec $(\ast)$ et utilisant à nouveau \unsure{a|à} la commutation de $\varinjlim_I$ à $\varinjlim_J$, on trouve donc
% manuscript position: 15
\[
g\big(\varinjlim\varphi(j)\big)\xleftarrow{\approx}\varinjlim g(\varphi(j)),
\]
% manuscript position: 16–22
ce qui \deleted{\unsure{achève de}} prouve que $g$ \deleted{est ds \unsure{$\underline{\mathrm{Hom}}$}} commute aux limites inductives de type $J$, et achève de prouver que $\unsure{O|g|0}$ est $\underline{\mathrm{Hom}}_{!}(\operatorname{Ind}_\pi(C),\unsure{\mathcal E|E})$. Comme $g\mid C\simeq f$, on \deleted{\unsure{en}} en conclut que le foncteur (4.4.16.2) est ess. surjectif, ce qui achève de prouver 4.4.16. De plus, on a construit un foncteur quasi-inverse explicite, $\psi$,
% manuscript position: foot
% archive sheet number: 188
\scanpage{190}
% manuscript position: top
\mspage{106}
% manuscript position: 1
s'insérant ds le diagramme commutatif
% manuscript position: 2–7
\diagram{{\scriptsize\begin{tikzcd}[ampersand replacement=\&,column sep=0.8em,row sep=large]
\underline{\mathrm{Hom}}_{!\pi}(C,\unsure{\mathcal E|E}) \arrow[r,"\psi","{(\approx)}"'] \arrow[d,hook,"\text{inc. pl. fid.}"] \& \underline{\mathrm{Hom}}_{!}(\operatorname{Ind}_\pi(C),\mathcal E)\arrow[r]\arrow[d,dashed,"\text{pl. fid}"] \& \underline{\mathrm{Hom}}_{!\pi}(\operatorname{Ind}_\pi(\unsure{C}),\mathcal E)\\
\underline{\mathrm{Hom}}(C,\mathcal E)\arrow[r,hook,"{\substack{L\\\approx}}","{\substack{\deleted{\text{pl. fid}}\\\deleted{\text{\unsure{d'image ess. formée|prolonge un foncteur}}}\\\deleted{\text{des foncteurs qui}}\\\deleted{\text{commute aux}}\\\deleted{\text{petites }\varinjlim\text{ filtrantes}}}}"'] \& \underline{\mathrm{Hom}}_{!\,\mathrm{filt}}(\operatorname{Ind}(\deleted{\unsure{C}}C),\mathcal E)\arrow[ur,"{R\ (\text{restriction})}"'] \& {}
\end{tikzcd}}}
% manuscript position: left/3
(4.4.16.\unsure{6|5})
% manuscript position: 8
[cf SGA 4 I 8.7.3]

% manuscript position: 9–10
4.5. \uline{Filtration cardinale canonique d'une} \inserted{catégorie} \inserted{$\mathcal U$-}\uline{paratopos}.

% manuscript position: 11–14
\uline{Définition} 4.5.1. On appelle \uline{\inserted{catégorie-$\mathcal U$-}para\deleted{\illegible{1}}topos} une catégorie \deleted{qui est} \inserted{$M$} \inserted{$\mathcal U$-}pseudo\deleted{\illegible{1}}topos et \deleted{\illegible{1}} accessible, i.e qui satisfait aux conditions équivalentes du théorème 4.3.\unsure{5|6} (p 70).
% manuscript position: 15–25
On dit \deleted{la description} \unsure{que|qu'}, c'est une \deleted{$\pi$-}catégorie $\pi$-\inserted{$\mathcal U$-}para\deleted{\illegible{1}}topos si elle est \inserted{\deleted{\unsure{pleine}}}$\pi$-accessible, i.e. si elle admet une petite sous-catégorie génératrice dont les objets sont $\pi$-\unsure{accessible|accessibles}, ou ce qui revient au même, si $M$ est équivalente à une catégorie de la forme $\operatorname{Ind}'_\pi(M_\pi)$, où $M_\pi$ est une \inserted{petite} catégorie [\inserted{ou : une cat.} équivalente à une petite catégorie], et stable par \deleted{petites} $\varinjlim$ de cardinal $\leq\pi$. (Et on peut prendre \deleted{\unsure{pour}} $M_\pi$ la sous-catégorie pleine de $M$ formée des objets
% manuscript position: foot
% archive sheet number: 189
\scanpage{191}
% manuscript position: top
\mspage{107}
% manuscript position: 1
$\pi$-\unsure{accessibles|accessible}

% manuscript position: 2–4
Dans la suite, nous allons, plus généralement, considérer une catégorie $C$ équivalente à une petite catégorie, \inserted{stable par $\varinjlim$ de cardinal $\leq\pi_0$,\unsure{$^{\ast}$}} et la catégorie
% manuscript position: left/2–5
% slanted
\margin{il n'y a pas lieu de le \unsure{supp.} \unsure{ça} tout de suite}
% manuscript position: 5
\[
M\overset{\mathrm{déf}}{=}\operatorname{Ind}'_{\pi_0}C\underset{\unsure{\neq|\not\simeq}}{\subset}\widehat C,\tag*{(4.5.1)}
\]
% manuscript position: 6
où $\pi_0$ est un cardinal infini donné.

% manuscript position: 7–8
Je me propose d'introduire sur $M$ une filtration cardinale canonique.

% manuscript position: 9–11
\uline{Proposition} 4.5.2. Soit $F\in\operatorname{Ob}M\subset\operatorname{Ob}\widehat C$, et soit $\pi$ un cardinal $\geq\pi_1=\operatorname{Sup}(\pi_0,\operatorname{card}\operatorname{Fl} C_0)$, \inserted{où $C_0$ est une \inserted{petite} sous-catégorie pleine \deleted{\unsure{réduite}} de $C$ telle que $C_0\overset{\approx}{\hookrightarrow}C$.} Les conditions suivantes sont équivalentes
% manuscript position: left/9–13
% slanted
\margin{N'est \uline{prouvé} (et n'est vrai) que si $\pi^{\pi_0}=\pi$ \unsure{4.5.5|4.5.6}. Voir rectification page 116}

% manuscript position: 12
\[
\text{a)}\quad\operatorname{Card}\operatorname{Fl}(C_{0/F})\leq\pi
\]

% manuscript position: 13
\[
\text{a$'$)}\quad\operatorname{Card}\deleted{\unsure{F}}\operatorname{Ob}(C_{0/F}^{\unsure{\circ|)}})\leq\pi
\]

% manuscript position: 14
a$''$) Pour tout $x$ dans $C_0$, $\operatorname{card} F(x)\leq\pi$

% manuscript position: 15–18
b) Il existe une petite catégorie filtrante $I$ (\unsure{par une|par ex.} ordonnée), \uline{grande devant $\pi_0$}, \inserted{avec \uline{$\operatorname{card} I\leq\pi$},} et un système inductif $(x_i)_{i\in I}$ dans $C$, tel que l'on ait
% manuscript position: 19
\[
\deleted{\unsure{x}}F\simeq\varinjlim_{i\in I}x_i,\tag*{(4.5.2.1)}
\]
% manuscript position: 20–24
\deleted{\unsure{dans}} la $\varinjlim$ est prise dans $M$, ou ce qui revient au \unsure{$=$} (en vertu de 4.4.15, $I$ étant gde devant $\pi_0$) dans $\widehat C$, i.e. $F$ est isomorphe au foncteur ind-représentable défini par le Ind-objet $\underline{x}=(x_i)_{i\in I}$ de $C$.

% manuscript position: 25
c) $F$ est $\pi$-accessible dans $M$

% manuscript position: 26
c$'$) $F$ est $\pi$-accessible dans $\widehat C$,
% manuscript position: foot
% archive sheet number: 190
\scanpage{192}
% manuscript position: top
\mspage{108}
% manuscript position: 1
\uline{Dém}. On a les implications $\pm$ tautologiques
% manuscript position: 2–3
\diagram{\begin{tikzcd}[ampersand replacement=\&,column sep=small,row sep=tiny]
\deleted{\unsure{a}}\arrow[r,phantom,"\deleted{\unsure{\Longrightarrow}}" description]\arrow[dd,phantom,"\deleted{\unsure{\Downarrow}}" description] \& a'\arrow[r,Rightarrow] \& a'' \& {}\\
{(\alpha)}\quad a\arrow[ur,Rightarrow]\arrow[dr,Rightarrow] \& {} \& {} \& {}\\
\deleted{b} \& b\arrow[r,Rightarrow] \& c'\arrow[r,Rightarrow] \& c^{\deleted{\unsure{\prime}}}\text{,}
\end{tikzcd}}
% manuscript position: 4–9
où $a'\Rightarrow a''$ provient de l'inclusion canonique $F(x)=\operatorname{Hom}_M(x,F)\mathrel{\lhook\joinrel\relbar}\operatorname{Ob} C_{/F}$, $a\Rightarrow b$ s'obtient en prenant $I=C_{0/F}$ et notant que $I$ est stable par limites inductives \deleted{grandes} \inserted{de} cardinal $\leq\pi_0$, ($C$ l'étant), donc est \unsure{grde|grdes} devant $\pi_0$.
% manuscript position: left/6–10
% slanted
\margin{on n'a pas besoin de cette stabilité, de \unsure{ttes} façons $C_{0/F}$ est \unsure{grde} devant $\pi_0$, \unsure{en vertu de|vu} 4.4.3 p. 80}

% manuscript position: 9–11
On a $b)\Rightarrow c'$ en vertu du fait que la sous-catégorie \inserted{str. pleine} $\unsure{\widehat M_\pi|\widehat C_\pi}$ de $\unsure{\widehat C|\widehat M}$ définie par
% manuscript position: 12–13
\[
\left\{\begin{aligned}
\widehat M_\pi={}&\text{sous-cat. \inserted{str.} pleine de }\widehat C\text{ formée des}\\
&F\in M\ \deleted{\text{\unsure{tels}}}\text{ qui sont }\pi\text{-accessibles}
\end{aligned}\right.\tag*{(\unsure{4.5.2.1|4.2.2.1})}
\]
% manuscript position: 14–19
est stable par $\varinjlim$ de cardinal $\leq\pi$ (SGA 4 I 9.9). Enfin $c'\Rightarrow c^{\unsure{\prime\prime|\prime}}$ provient du fait que le \deleted{\unsure{lim}} foncteur d'inclusion $M\hookrightarrow\widehat C$ est pl. fidèle et \uline{$\pi_0$-accessible}, i.e. commute aux limites inductives grandes devant $\pi_0$, et a fortiori à celles grandes devant $\pi$.

% manuscript position: 20–22
\deleted{Donc pour \unsure{achever}} \deleted{\unsure{la}} Par les indications du diagramme ci-dessus, il suffit donc de prouver qu'on a
% manuscript position: left/22
\margin{$\unsure{2|Z}$}
% manuscript position: 23
\[
c\overset{\inserted{\text{n'est prouvé (et n'est vrai) que si }\pi^{\pi_0}=\pi}}{\Rightarrow} a
\qquad\text{et}\qquad a''\Rightarrow a.
\]
% manuscript position: 24–26
\deleted{On \unsure{va|en} d'abord \unsure{prouver|prouvons} \deleted{\unsure{$a\Rightarrow$}} $a''\Rightarrow a$, ce qui implique déjà \unsure{l'équivalence|l'égalisation} des}
% manuscript position: foot
% archive sheet number: 191
\scanpage{193}
% manuscript position: top
\mspage{109}
% manuscript position: 1–3
\deleted{des cinq conditions, autres que $c$, soit $P(\pi)$ la condition obtenue et l'implication $P(\pi)\Rightarrow c$. Il restera à prouver encore}
% manuscript position: 4
\[
\deleted{c\Rightarrow P(\pi).}
\]

% manuscript position: 5
\uline{Prouvons} $a''\Rightarrow a$. Considérons l'application
% manuscript position: 6
\[
\operatorname{Fl} C_{/F}\longrightarrow\operatorname{Fl} C\tag*{(\unsure{\(\alpha\)|\(\ast\)})}
\]
% manuscript position: 7–16
induite par le foncteur \inserted{étale} $C'=C_{/F}\longrightarrow C$, alors l'ens. des flèches au dessus d'une flèche $\alpha:x\longrightarrow y$ de $C$ est \deleted{de cardinal} en corr. biunivoque avec $\operatorname{Ob} C'_y=F(y)$, puisqu'une telle flèche $\alpha'$ est connue quand on connaît son but, qui est un objet quelconque de $F(y)$ ($C'\longrightarrow C$ étant \unsure{étale} i.e. fibré à fibres discrètes). Donc la fibre de \unsure{(\(\alpha\))|(\(\ast\))} en $\alpha$ est \deleted{majorée} de cardinal \deleted{$\leq$} $\operatorname{card} F(y)\leq\pi$, donc on a
% manuscript position: 17
\[
\operatorname{Card}\operatorname{Fl} C_{/F}\leq\pi\underbrace{\operatorname{card}\operatorname{Fl} C}_{\leq\pi_1}
\leq\pi\underbrace{\pi_1}_{\leq\pi}\leq\pi^2=\pi,
\]
% manuscript position: 18
ce qui prouve $a$.

% manuscript position: 19–21
Prouvons enfin $c\Rightarrow a$. Pour ceci, nous allons utiliser \inserted{la première assertion \unsure{du n°|du no} $M_\pi\subset\operatorname{Filt}^{\unsure{\geq\pi|\pi}} M$ de} SGA 4 I 9.15, \deleted{qui} s'applique dès lors que nous prouvons le
% manuscript position: left/22–27
% slanted
\margin{\unsure{2|Z} Pas prouvé ; la condition c) \unsure{aux filtrations|aux filtrations sur} \unsure{cardinales|cardinaux} paraît-être \unsure{pas satisfaite|en défaut}. \deleted{\unsure{Sans doute}} Faut \unsure{affiner|rectifier} \deleted{\illegible{1}}}

% manuscript position: 22–27
\uline{Corollaire} 4.5.3. Sous les conditions de 4.5.2, soit \deleted{\unsure{,}} pour tout cardinal $\pi\geq\pi_1$, $\deleted{\unsure{M_\pi}}\operatorname{Filt}^{\pi}M$ la sous-\uline{\deleted{catégorie strictement pleine}} \inserted{\deleted{\unsure{ensemble}}} de $\operatorname{Ob} M$ formé des $F\in\operatorname{Ob}M$ \deleted{\unsure{$\pi$-accessibles (une des conditions)}} qui satisfont la condition a) de 4.5.2, Alors la famille des $(\operatorname{Filt}^{\pi}M)_{\pi\geq\pi_1}$ est une
% manuscript position: foot
% archive sheet number: 192
\scanpage{194}
% manuscript position: top
\mspage{110}
% manuscript position: 1
filtration cardinale de $M$.

% manuscript position: 2–4
La condition a) de loc. cit. 9.12 (les $\underline{\mathrm{Filt}}^\pi$ sont équivalentes à des petites catégories) est évidente. Prouvons \inserted{b), i.e.} la stabilité \unsure{de|par} $\underline{\mathrm{Filt}}^\pi$ par $\varinjlim$ filtrantes \inserted{$F=\varinjlim F_i$} \uline{grdes devant $\pi_0$}, de cardinal $\leq\pi$.
% manuscript position: 5–9
Comme $M\subset \widehat C$ \deleted{est} commute à ces limites, \deleted{il suffit de voir que $\operatorname{Card} F(x)$} pour $x$ de $C$
% manuscript position: 10
\[ F(x)=\varinjlim_{\unsure{i|x}} F_i(x) \]
% manuscript position: 11–13
d'où $\operatorname{card}F(x)\leq\operatorname{card}\operatorname{Ob} I\cdot\pi\leq\pi\cdot\pi=\pi$, ce qui (par l'équivalence $a\Longleftrightarrow a''$) implique $F\in\operatorname{Ob}\mathrm{Filt}^\pi$.

% manuscript position: 14–17
Prouvons la condition c) \inserted{de loc. cit.}. Pour ceci, nous avons besoin du fait, \unsure{$(^{\ast1})$} plus précis que le b) qu'on vient d'établir :
% manuscript position: left/18–20
% slanted
\margin{Prouvé seulement si $\pi^{\pi_0}=\pi$ ! Sans doute Faux en général.}

% manuscript position: 18–21
[ \uline{Corollaire} 4.5.4. \deleted{\unsure{Pour|Pour $\mathrm{Filt}^\pi$}} \deleted{$\unsure{\mathrm{Filt}^\pi}$} tout $\pi\geq\pi_1$, $\mathrm{Filt}^\pi$ \inserted{(dans $M$)} est stable par $\varinjlim$ de cardinal $\leq\pi$ (pas néc. filtrantes, a fortiori pas néc. grandes devant $\pi_0$).
% manuscript position: left/21–25
% slanted
\margin{Il faut une variante des \unsure{énoncés|critères} 4.4.14, pour \unsure{qui ne présuppose|qui se prouve} \unsure{pas que} $C$ \unsure{n'ait pas G.S.|n'est pas G.S.} ?}

% manuscript position: 22–24
Il suffit, \deleted{\unsure{de $+$|de}} en vertu de ce qui est déjà prouvé et du lemme $4.4.14^{\ast}$, de montrer que $\mathrm{Filt}^\pi$ est stable par $\varinjlim$ \uline{de cardinal} $\leq\pi_0$.
% manuscript position: 25–27
Si $\varphi:J\longrightarrow\mathrm{Filt}^{\unsure{=\pi|\pi}}$ est un foncteur avec $\operatorname{card}J\leq\pi_0$ on le met sous forme indicielle
% manuscript position: 28
\[ \psi:I\times J\longrightarrow C \]
% manuscript position: 29
grâce à \unsure{4.4.13} b), avec $I$ grd devant $\pi_0$.
% manuscript position: foot
% archive sheet number: 193
\scanpage{195}
% manuscript position: top
\mspage{111}
% manuscript position: 1
Je rappelle qu'on pouvait prendre
% manuscript position: 2
\[ I=\underline{\mathrm{Hom}}(J,C)_{/\varphi}\,, \]
% manuscript position: 3–4
où $\underline{\mathrm{Hom}}(J,C)\lhook\joinrel\longrightarrow\underline{\mathrm{Hom}}(J,\text{\deleted{$\unsure{\operatorname{Ind}_{\pi}|C}$}}\,M\deleted{\unsure{)\longrightarrow|)}}=\operatorname{Ind}'_{\pi_0}(C))$, quitte à \deleted{\unsure{ne}} pas exiger que $I$ soit ordonnée.

% manuscript position: 5–7
On a envie de prouver que $\operatorname{card}\mathrm{Fl}\,I\leq\pi$. L'ennui, c'est qu'on arrive seulement à la majoration
% manuscript position: 8
\[ \operatorname{card}\mathrm{Fl}\,I\leq\pi^{\pi_0}. \]
% manuscript position: 9–10
\deleted{(moyennant l'hypothèse $\pi\geq\pi_1^{\unsure{\pi_0|\pi}}$, \unsure{de f.|de fait} (c'est \unsure{(}} \deleted{\unsure{M}}
% manuscript position: 10–11
si $J$ est discrète, on n'a pas mieux, puisqu'alors
% manuscript position: 12
\[ I\simeq\prod_{j\in J} C_{/\varphi(j)}
\qquad\text{d'où }\mathrm{Fl}\,I=\prod_{j\in\unsure{I|J}}\mathrm{Fl}\,C_{/\varphi(j)}, \]
% manuscript position: 13–16
et si $\operatorname{Card}\mathrm{Fl}\,C_{/\varphi(j)}=\pi$ pour $\unsure{\forall|\mathrm{tt}}\unsure{,}\,j$, on trouve $\operatorname{card}\mathrm{Fl}\,I=\pi^{\pi_0}$ -- il n'y a pas moyen de faire mieux ! ]

% manuscript position: 17–19
Donc je n'arrive pas à prouver directement 4.5.4, ce qui m'aurait arrangé pour prouver la condition c) de loc. cit.
% manuscript position: 20–25
Mais, à vrai dire, \unsure{en|on} \deleted{\unsure{a vu déjà|a}} \unsure{rien que|déjà|une fois} par le travail de la démonstration du th. 4.3.6, partie $b\Rightarrow c$ -- on y a déjà prouvé que la filtration qu'on vient de construire là (par la condition a) \deleted{\unsure{que}} de 4.5.2), est une filtration cardinale, \deleted{La} la condition c) de SGA 4 I 9.12 est prouvée \unsure{à} la page 76. Mais je vais refaire
% manuscript position: foot
% archive sheet number: 194
\scanpage{196}
% manuscript position: top
\mspage{112}
% manuscript position: 1
la démonstration plus \unsure{sérieusement|sérieusemt}.

% manuscript position: 2
Soit donc $F$ \deleted{$\unsure{\in\mathrm{Filt}^\pi|=\mathrm{Filt}^\pi}$} \inserted{$\in\operatorname{Ob} M$}, \deleted{\unsure{considérons}} considérons
% manuscript position: 3
\[ J=C_{/F}\,,\quad \text{\deleted{($\unsure{\operatorname{card}\mathrm{Fl}\,J\leq\pi|{}\text{de }\operatorname{card}\mathrm{Fl}\,J\leq\pi}$)}}, \]
% manuscript position: 4–5
c'est une catégorie stable par \deleted{\unsure{sommes directes}} \inserted{$\varinjlim$ de cardinal} $\leq\pi_0$, donc grande devant $\pi_0$, et on a
% manuscript position: 6
\[ F\simeq\varinjlim_{j\in J}x_j\quad(\varinjlim\text{ dans }M) \]
% manuscript position: 7–8
en désignant par $j\longmapsto x_j$ le foncteur canonique $C_{/F}\longrightarrow C$.

% manuscript position: 9–12
Soit donc $I$ l'ens des sous-catégories \deleted{\unsure{pleines}}\inserted{$J'$} de $J$ qui sont \uline{de cardinal} $\leq\pi$ \uline{et} \uline{grandes devant} $\pi_0$, $I$ ordonné par inclusion, et pour tout $J'\in I$, posons
% manuscript position: 13
\[ X_{J'}=\varinjlim_{j\in J'}x_j\,. \]
% manuscript position: 14–16
\deleted{On a alors} \deleted{Alors \unsure{$\forall$} $J'\longmapsto X_{J'}$ est un foncteur $I\longrightarrow\mathrm{Filt}$} \deleted{\unsure{Cette|Celle}}

% manuscript position: 16–20
Comme $J'$ est grand devant $\pi_0$, la limite inductive dans $M$ est aussi une limite dans $\widehat C$, d'où résulte \deleted{que} aussitôt, comme $\operatorname{card}\,\text{\deleted{$\unsure{F|J}$}}\,\operatorname{Ob} J'\mathrel{\unsure{\subset|<}}\operatorname{card}\mathrm{Fl}\,J'\leq\pi$, que $X_{J'}$ est dans $\mathrm{Filt}^\pi$. Alors
% manuscript position: 21
\[ J'\longmapsto X_{J'}\in\operatorname{Ob}\mathrm{Filt}^{\unsure{\pi|n}} \]
% manuscript position: 22
est un foncteur
% manuscript position: 23
\[ I\longrightarrow\mathrm{Filt}^\pi\,, \]
% manuscript position: 24–27
dont la limite inductive est celle des $x_j$, donc isom. \unsure{;|à} $F$. D'autre part il est immédiat que $I$ est grd devant $\pi$, i.e. que si $(J'_\alpha)_{\alpha\in A}$ \deleted{est}
% manuscript position: foot
% archive sheet number: 195
